\chapter{Spectral Functions, Kubo Formulae, and Transport}
\label{ch:thermal-spectral-kubo-transport}

\ConventionsRealTime

Transport coefficients are low-frequency limits of real-time response
functions.  The definitions use a thermal state, a conserved current or
stress tensor, and a source coupled to the corresponding operator.  The
derivation is linear response from the Schwinger--Keldysh generating
functional of Chapter~\ref{ch:thermal-schwinger-keldysh}.

\section{Retarded Functions and Spectral Densities}

Let \(\rho_\beta\) be a KMS state at inverse temperature \(\beta\).  For
bosonic Hermitian operators \(A(t,\mathbf x)\) and \(B(t,\mathbf x)\), define
the retarded function
\[
  G^R_{AB}(t,\mathbf x)
  =
  -\ii\,\Theta(t)\,
  \langle [A(t,\mathbf x),B(0,\mathbf 0)]\rangle_\beta
\]
and the spectral density
\[
  \rho_{AB}(t,\mathbf x)
  =
  \langle [A(t,\mathbf x),B(0,\mathbf 0)]\rangle_\beta .
\]
Fourier transforms are defined distributionally by
\[
  G^R_{AB}(\omega,\mathbf k)
  =
  \int dt\,d^{D-1}x\,
  e^{\ii\omega t-\ii\mathbf k\cdot\mathbf x}
  G^R_{AB}(t,\mathbf x).
\]
For Hermitian \(A=B\), the finite-volume spectral representation gives
\[
  \rho_{AA}(\omega,\mathbf k)
  =
  \frac{2\pi}{Z}\sum_{m,n}
  e^{-\beta E_m}
  \left(1-e^{-\beta\omega}\right)
  |\langle m|A_{\mathbf k}|n\rangle|^2
  \delta(\omega-E_n+E_m),
\]
so \(\omega\,\rho_{AA}(\omega,\mathbf k)\geq0\) as a distributional
positivity statement after the thermodynamic limit is taken under the stated
hypotheses.

\section{KMS and Fluctuation--Dissipation}

Define the Wightman functions
\[
  G^>_{AB}(t)=\langle A(t)B(0)\rangle_\beta,\qquad
  G^<_{AB}(t)=\langle B(0)A(t)\rangle_\beta .
\]
The KMS condition gives
\[
  G^>_{AB}(t-\ii\beta)=G^<_{AB}(t).
\]
In frequency space this implies
\[
  G^>_{AB}(\omega)=e^{\beta\omega}G^<_{AB}(\omega),
  \qquad
  \rho_{AB}(\omega)=G^>_{AB}(\omega)-G^<_{AB}(\omega).
\]
Therefore the symmetrized correlator
\[
  G^{\rm sym}_{AB}(\omega)
  =
  \frac12\bigl(G^>_{AB}(\omega)+G^<_{AB}(\omega)\bigr)
\]
satisfies
\[
  G^{\rm sym}_{AB}(\omega)
  =
  \frac12\coth\!\left(\frac{\beta\omega}{2}\right)\rho_{AB}(\omega).
\]

\section{Linear Response}

Perturb the Hamiltonian by a source \(h_B(t,\mathbf x)\) coupled to \(B\):
\[
  H(t)=H_0-\int d^{D-1}x\,h_B(t,\mathbf x)B(t,\mathbf x).
\]
To first order in \(h_B\), the change in the expectation value of \(A\) is
\[
  \delta\langle A(t,\mathbf x)\rangle
  =
  \int dt' d^{D-1}x'\,
  G^R_{AB}(t-t',\mathbf x-\mathbf x')h_B(t',\mathbf x').
\]
This follows by expanding the time-evolution operator to first order on the
Schwinger--Keldysh contour and subtracting the insertion on the backward
branch.  The contour difference is exactly the commutator, and causality is
the step function.

\section{Conductivity and Viscosity}

For a conserved current \(J^\mu\) coupled to a background gauge field
\(A_\mu^{\rm ext}\), the electric field is
\[
  E_i(\omega,\mathbf k)=\ii\omega A_i^{\rm ext}(\omega,\mathbf k)
  -\ii k_i A_0^{\rm ext}(\omega,\mathbf k).
\]
At zero spatial momentum and in a gauge with \(A_0^{\rm ext}=0\), linear
response gives
\[
  \langle J_i(\omega)\rangle
  =
  G^R_{J_iJ_j}(\omega,\mathbf 0)A_j^{\rm ext}(\omega)
  =
  \sigma_{ij}(\omega)E_j(\omega),
\]
hence
\[
  \sigma_{ij}(\omega)
  =
  \frac{1}{\ii\omega}G^R_{J_iJ_j}(\omega,\mathbf 0)
\]
after subtracting possible contact terms dictated by gauge invariance.

For the stress tensor, couple the theory to a spatial metric perturbation
\(g_{ij}=\delta_{ij}+h_{ij}\).  In an isotropic state, the shear viscosity is
the zero-frequency slope
\[
  \eta
  =
  \lim_{\omega\to0^+}
  \frac{1}{\omega}\,
  \operatorname{Im}G^R_{T_{xy}T_{xy}}(\omega,\mathbf 0).
\]
The bulk viscosity is obtained from the traceful spatial stress channel after
projecting out pressure contact terms and conserved-density mixing.

\section{Hydrodynamic Pole Boundary}

Kubo formulae define transport coefficients when the indicated limits exist
and commute with the thermodynamic and zero-momentum limits in the stated
order.  Hydrodynamics begins after these coefficients are inserted into the
long-wavelength Ward identities for conserved densities.

\begin{openproblem}[Rigorous Kubo-to-hydrodynamics bridge]
\label{op:kubo-hydrodynamics-bridge}
Give sufficient QFT hypotheses under which retarded stress-tensor and current
correlators have the analyticity, spectral positivity, clustering, and
low-frequency limits required by the Kubo formulae, and prove that the same
coefficients control the long-time hydrodynamic effective theory.
\end{openproblem}
